\section{Floating-Point Representation and Numerical Errors}\label{sec:floating-point}

Computers approximate real numbers using finite-precision representations 
governed by the IEEE 754 standard for floating-point arithmetic\cite{IEEE754_2019}.
Compared to integer representations, floating-point numbers can represent a much wider range of values,
by using a normalized scientific notation with a fixed number of significant digits (the mantissa) and an exponent.

Floating-point systems introduce fundamental error sources\cite{Goldberg1991FloatingPoint}:
\begin{itemize}
\item \textbf{Rounding and Machine Epsilon:} 
    Real numbers cannot be represented with infinite precision. 
    The spacing between adjacent representable floating-point numbers is bounded by the Machine Epsilon $\epsilon_{\text{mach}}$.
    Operations on physical quantities introduce rounding errors that accumulate over multistep time integrations 
    (e.g., Verlet or Runge-Kutta integrators).
    While single precision (32-bit) is the de facto standard for most programming languages, 
    it is insufficient for large-scale simulations with both high number ranges and small deltas.
    Double precision (64-bit) and quadruple precision (128-bit) are also available on modern processors, 
    but come with increased memory usage and computational cost. 
    See \Cref{tab:floating_point_comparison} for a comparison of IEEE 754 formats.
\item \textbf{Catastrophic Cancellation:}
    Subtracting two nearly equal floating-point numbers results in loss of significant digits. 
    This frequently occurs when computing relative position vectors between close particles 
    ($\vec{r}_1 - \vec{r}_2$), leading to increased relative errors.
\item \textbf{Underflow, Overflow, and NaN:}
    Exceeding the dynamic range of the exponent field results in overflow (yielding $\pm\infty$) 
    or underflow (yielding subnormal numbers or zero), which can corrupt calculations.
    A floating-point values might also hold the special value NaN (Not a Number) to indicate undefined or unrepresentable results,
    such as $0/0$ or $\infty - \infty$.
\end{itemize}

\begin{table*}
\centering
\small
\caption{Comparison of IEEE 754 floating-point formats\cite{IEEE754_2019}}
\label{tab:floating_point_comparison}
\begin{tabularx}{\textwidth}{l c c X X}
\toprule
\textbf{Format} & \textbf{Digits} & \textbf{($\epsilon_{\text{mach}}$)} & \textbf{Dynamic Range} & \textbf{Relative Error ($50^{\text{th}}$ / $99^{\text{th}}$)} \\
\midrule
\texttt{f32}  & $6-7$   & $\approx 1.19 \times 10^{-7}$  & $1.18 \times 10^{-38}$   to $3.40 \times 10^{38}$   & $\sim 10^{-7}$  / $\sim 10^{-5}$  \\
\texttt{f64}  & $15-17$ & $\approx 2.22 \times 10^{-16}$ & $2.23 \times 10^{-308}$  to $1.79 \times 10^{308}$  & $\sim 10^{-16}$ / $\sim 10^{-14}$ \\
\texttt{f128} & $34$    & $\approx 1.92 \times 10^{-34}$ & $3.36 \times 10^{-4932}$ to $1.18 \times 10^{4932}$ & $\sim 10^{-34}$ / $\sim 10^{-32}$ \\
\bottomrule
\end{tabularx}
\end{table*}

